% !TeX root = ../../main.tex
\subsection{方程式の利用の問題について}
方程式の利用の問題パターンはせいぜい， \fitblankbf[5パターン]{} 程度である.
それは，

\begin{enumerate}
    \item \fitblankbf[個数と代金]{} の問題
    \item \fitblankbf[道のり・速さ・時間]{} の問題
    \item \fitblankbf[割合]{} の問題
    \item \fitblankbf[過不足]{} の問題
    \item \fitblankbf[整数]{} の問題
\end{enumerate}

ここで，それぞれの問題について深掘りする前に，前提となる考え方を共有しておこう.
\vspace{2em}

\begin{techniquebox}[tech:01]{\textbf{方程式は「\textbf{天秤}」というイメージを持とう！}}
    例えば，1 + 1 = 2 という式があったときに，（途中式）=（答え）というように，
    左から右の一方通行をイメージしていないだろうか.

    しかし，「=」という記号は数学的に，\textbf{左辺にある数量}と\textbf{右辺にある数量}
    とが見た目は違えど，中身が\textbf{等しい}という意味を持つ.

    したがって，方程式は「\textbf{天秤}」というイメージを持ち，重い方から少し取ってバランスを取ったり，
    軽い方に少し足してバランスを取ったりするという考え方を持つことで立式しやすくなる.
\end{techniquebox}
\newpage
